\scanpage{10}\mspage{9}
% ink top, lines 1–3
a \uline{half}-\uline{interval} if it admits a least element but not a greatest, or else the reverse.

% ink upper middle, NB, lines 1–4
\deleted{\unsure{The opposite|The opposite}}\inserted{\uline{NB}} The opposite of a segment piece, segment, interval, half-interval is likewise. The opposite of a $1/2$ interval closed on the left is closed on the right, and conversely.
% ink left, slanting upwards beside NB
% slanted
\margin{In my current terminology, segment piece \unsure{is|to be|was|to be} \unsure{``set|set|a|équiv} \uline{biorder}}

% ink middle, Tout, lines 1–3
\deleted{\unsure{When|When}}\inserted{Every} segment piece is one of the three (segment, $1/2$ interval, interval) to the exclusion of the other two.

% ink lower middle, Si L
If $\mathcal L$ is a segment piece, we put $\deleted{\unsure{\partial L}=}$
\[
\begin{aligned}
\unsure{\partial|\partial L}&=\bigl\{\,x\in L\mid
\begin{aligned}&x\text{ is least element}\\
&\text{or greatest element}\end{aligned}\,\bigr\}
\end{aligned}
\]
\deleted{\illegible{1}}

% ink bottom, Soit and first equality
Let $\mathcal L$ be a segment piece, and
\[
\begin{aligned}
S_0(\mathcal L)=\mathrm{Drap}_2(\mathcal L)&=\{\,\unsure{\varepsilon|\mathcal E|E}\in P_2(\mathcal L)\mid \unsure{\varepsilon|\mathcal E|E}\text{ totally ordered}\,\}\\
&=\{\,\{x,y\}\mid x<y\,\}
\end{aligned}
\]
% ink bottom, final equality
\[
\text{Whence}\qquad \mathrm{Modf}(\mathcal L)=\mathrm{Drap}(\mathcal L)
=\{\,S\subset P_f(\mathcal L)\mid S\,\substack{\text{totally}\\\text{ordered}}\,\}
\]

% pale vertical paper watermark, rotated for reading
% the paper's watermark (EXTRA STRONG), not his text
